% 英文版特有的设置：章名、书名、题注名、定理名、各章引言、页脚、封面。由 main-en.tex 在 style/book.tex 之后 \input。
% 中文版的对应设置在 style/lang-zh.tex，两个文件里的条目一一对应，改了一个，记得看看另一个。

\title{C Language: A Lazy Approach}
\hypersetup{pdftitle={C Language: A Lazy Approach}, pdfauthor={whensungoesdown}}

% 文档类仍是 ctexbook（全书的版式、中文字体都靠它），这里把它的“第一章”“目录”“图”“表”等换成英文
\ctexset{
  contentsname   = Contents,
  listfigurename = List of Figures,
  listtablename  = List of Tables,
  figurename     = Figure,
  tablename      = Table,
  bibname        = Bibliography,
  indexname      = Index,
  appendixname   = Appendix,
  chapter  = {name = {Chapter~,},  number = \arabic{chapter}},
  appendix = {name = {Appendix~,}, number = \Alph{chapter}},
}

\SetupFloatingEnvironment{listing}{name=Listing, listname=List of Listings, within=chapter}

% 定理环境，按章编号
\newtheorem{theorem}{Theorem}[chapter]
\newtheorem{definition}[theorem]{Definition}
\newtheorem{lemma}[theorem]{Lemma}
\renewcommand{\proofname}{Proof}

% 页脚：偶数页书名，奇数页作者，章首页只写书名
\newcommand\lazyfooteven{C Language: A Lazy Approach}
\newcommand\lazyfooterodd{whensungoesdown}
\newcommand\lazyfooteropener{C Language: A Lazy Approach}

% 封面上不印中文书名（frontmatter/cover.tex 里用 \providecommand 定义，这里先定义为空）
\newcommand\CoverZhTitle{}

% 每章章首页的引言：\lazyquote{章号}{}{原文}{出处}，第二项（中文）留空
\lazyquote{1}{}{The only way to learn a new programming language is by writing programs in it.}
             {Brian W. Kernighan and Dennis M. Ritchie,\\\emph{The C Programming Language}, 1978}
\lazyquote{2}{}{Programs must be written for people to read, and only incidentally for machines to execute.}
             {Harold Abelson and Gerald Jay Sussman,\\\emph{Structure and Interpretation of Computer Programs}, 1985}
\lazyquote{3}{}{Everyone knows that debugging is twice as hard as writing a program in the first place.
              So if you're as clever as you can be when you write it, how will you ever debug it?}
             {Brian W. Kernighan and P. J. Plauger,\\\emph{The Elements of Programming Style}, 2nd ed., 1978}
\lazyquote{4}{}{All problems in computer science can be solved by another level of indirection.}
             {David J. Wheeler\\(quoted by Butler Lampson in his Turing Award lecture)}
\lazyquote{5}{}{This is the Unix philosophy: Write programs that do one thing and do it well.
              Write programs to work together.}
             {Doug McIlroy, in Peter H. Salus,\\\emph{A Quarter Century of Unix}, 1994}
\lazyquote{A}{}{I'm doing a (free) operating system (just a hobby, won't be big and professional like gnu)
              for 386(486) AT clones.}
             {Linus Torvalds, newsgroup comp.os.minix,\\August 25, 1991}
\lazyquote{B}{}{This is an attempt to stop a war. I hope it is not too late and that somehow,
              magically perhaps, peace will prevail again.}
             {Danny Cohen, \emph{On Holy Wars and a Plea for Peace},\\April 1, 1980}
\lazyquote{C}{}{Any fool can write code that a computer can understand.
              Good programmers write code that humans can understand.}
             {Martin Fowler,\\\emph{Refactoring}, 1999}
